\documentclass[10pt]{article}
\usepackage[T1]{fontenc}
\usepackage[utf8]{inputenc}
\usepackage[french]{babel}
\usepackage{amsmath,amssymb,mathrsfs,mathpazo}
\usepackage{geometry}
\geometry{a4paper,margin=2cm}
\usepackage{graphicx,array,tabularx,tabularray,colortbl,multirow,multicol,booktabs,enumitem,pifont,ulem,fancybox,diagbox,cancel}
\usepackage{pstricks,pstricks-add,pst-plot,pst-tree,pst-node,pst-text,pst-func,pst-math,pst-bspline,pst-3dplot,pst-eucl,pst-solides3d}
\let\pscalculate\undefined
\usepackage{pst-bezier}
\usepackage{tikz,pgfplots,tkz-tab}
\usetikzlibrary{arrows,arrows.meta,patterns,calc,positioning,trees,decorations.markings,shapes.geometric,shapes.misc,shapes.arrows,intersections,decorations.pathmorphing,decorations.pathreplacing,backgrounds,shadows,babel}
\pgfplotsset{compat=1.18}
\usepgfplotslibrary{fillbetween}
\usepackage[np]{numprint}
\usepackage{listings,makecell,pas-tableur,twoopt}
\let\barre\undefined
\usepackage[pstricks]{tablvar}
\definecolor{darkgreen}{RGB}{0,100,0}
\providecommand{\red}{\color{red}}
\providecommand{\blue}{\color{blue}}
\providecommand{\green}{\color{green}}
\providecommand{\black}{\color{black}}
\providecommand{\white}{\color{white}}
\providecommand{\gray}{\color{gray}}
\providecommand{\cyan}{\color{cyan}}
\providecommand{\magenta}{\color{magenta}}
\providecommand{\yellow}{\color{yellow}}
\providecommand{\violet}{\color{violet}}
\providecommand{\orange}{\color{orange}}
\providecommand{\euro}{\texteuro}
\providecommand{\EUR}{\texteuro}
\providecommand{\ieme}{\textsuperscript{e}}
\providecommand{\degres}{\ensuremath{^\circ}}
\providecommand{\degree}{\ensuremath{^\circ}}
\providecommand{\vv}[1]{\overrightarrow{#1}}
\providecommand{\Coord}[1]{#1}
\providecommand{\decosix}{}
\providecommand{\starredbullet}{\ensuremath{\bullet}}
\pagestyle{empty}
\setlength{\parindent}{0pt}
\usepackage{adjustbox}\DeclareUnicodeCharacter{2236}{:}\DeclareUnicodeCharacter{2190}{\ensuremath{\leftarrow}}
\providecommand{\R}{\mathbb{R}}
\providecommand{\N}{\mathbb{N}}
\providecommand{\Z}{\mathbb{Z}}
\providecommand{\Q}{\mathbb{Q}}
\providecommand{\D}{\mathbb{D}}
\providecommand{\C}{\mathbb{C}}
\providecommand{\e}{{\mathrm{e}}}
\providecommand{\vect}[1]{\overrightarrow{#1}}
\providecommand{\vectt}[1]{\overrightarrow{\mathrm{#1}}}
\providecommand{\vvt}[1]{\overrightarrow{\mathrm{#1}}}
\providecommand{\np}[2][]{{#2~#1}}
\providecommand{\vv}[1]{{\overrightarrow{#1}}}
\renewcommand{\d}{{\mathrm{d}}}
\providecommand{\degres}{{^\circ}}
\providecommand{\degre}{{^\circ}}
\providecommand{\blue}{\color{blue}}
\providecommand{\red}{\color{red}}
\providecommand{\enskip}{\quad}
\providecommand{\llbracket}{\lbrack\!\lbrack}
\providecommand{\rrbracket}{\rbrack\!\rbrack}
\providecommand{\euro}{{\text{€}}}
\providecommand{\no}{{\text{n°}}}
\providecommand{\interval}[2]{{[#1;#2]}}
\providecommand{\intervFO}[2]{{[#1;#2[}}
\providecommand{\dotfill}{\ldots}

\providecommand{\ds}{\displaystyle}
\providecommand{\textding}[1]{#1}
\providecommand{\Oij}{\(\mathrm{O};\vec{\imath},\vec{\jmath}\)}
\providecommand{\Oijk}{\(\mathrm{O};\vec{\imath},\vec{\jmath},\vec{k}\)}
\providecommand{\Ouv}{\(\mathrm{O};\vec{u},\vec{v}\)}
\providecommand{\intervFF}[2]{\ensuremath{[#1;#2]}}
\renewcommand{\euro}{\ifmmode\text{\texteuro}\else\texteuro\fi}
\DeclareUnicodeCharacter{00B0}{\ifmmode^{\circ}\else\textdegree\fi}
\DeclareUnicodeCharacter{20AC}{\ifmmode\text{\texteuro}\else\texteuro\fi}
% Énoncés : Éducation nationale. Transcription et corrigés : APMEP, Denis Vergès et contributeurs. Découpage et classement : CoopMaths, Sébastien Lozano, Éric Elter et contributeurs. Adaptation de la présentation : Maths Explorer.
% CC BY-SA : https://coopmaths.fr/www/about/
% https://coopmaths.fr/alea/static/e3c/2020/tex/e3c_2020_00_sujet40_4.tex
\begin{document}
\noindent\textit{APMEP / CoopMaths. CC BY-SA. Présentation adaptée par Maths Explorer.}\par\medskip
\medskip

\begin{enumerate}
\item Soit $f$ la fonction définie sur $\R$ par $f(x) = 4x^3 - 48x^2 + 144x.$ 
	\begin{enumerate}
		\item Calculer $f'(x)$ et montrer que $f'(x) = 12\left(x^2 - 8x + 12\right)$.
		\item En déduire le tableau de variations de la fonction $f$ sur $\R$.
	\end{enumerate}
\item Dans une plaque de carton carrée de $12$~cm de côté, on découpe, aux quatre coins, des carrés identiques afin de construire une boîte sans couvercle, comme indiqué sur les figures ci-dessous.

On note $x$ la longueur (en cm) du côté de chacun des carrés découpés.

On admet que $x \in ]0~;~6[$.

\begin{center}
\psset{unit=0.3cm}
\begin{pspicture}(30,16)

\psframe(1,1)(13,13)
\pspolygon[linestyle=dashed,fillstyle=solid,fillcolor=lightgray](3,1)(3,3)(1,3)(1,11)(3,11)(3,13)(11,13)(11,11)(13,11)(13,3)(11,3)(11,1)
\psline{<->}(11,0.5)(13,0.5)\uput[d](12,0.5){$x$}
\psline{<->}(13.4,1)(13.4,3)\uput[r](13.4,2){$x$}
\psline{<->}(0.3,1)(0.3,13)\psline{<->}(1,13.6)(13,13.6)
\uput[l](0.3,7){12}\uput[u](7,13.6){12}
\psframe[fillstyle=solid,fillcolor=lightgray](17,1)(25,3)
\pspolygon[fillstyle=solid,fillcolor=lightgray](25,1)(30,3)(30,5)(25,3)
\psline(17,3)(22,5)(30,5)

\psline(22,5)(22,3)
\psarc[linewidth=0.15cm]{<-}(15.5,0){5}{70}{110}
\end{pspicture}
\end{center}


L'objectif est de déterminer la longueur $x$ permettant d'obtenir une boîte de volume maximal.
	\begin{enumerate}
		\item Montrer que le volume de la boîte est égal à $100$ cm$^3$ pour $x = 1$. Détailler le calcul.
		\item Montrer que, pour $x \in ]0~;~6[$, le volume de la boîte est égal à $f(x)$,\: $f$ étant la fonction étudiée à la question 1.
		\item Quelle est la valeur de $x$ permettant d'obtenir une boîte de volume maximal ?
	\end{enumerate}
\end{enumerate}

\end{document}
